% Reverse-mode differentiation of the TF24 size- and trait-structured model.
% Written 2026-08-04. NOT COMPILED: no TeX toolchain in the authoring container,
% so this source has not been typeset and may carry syntax errors.
\documentclass[11pt,a4paper]{article}
\usepackage[margin=1in]{geometry}
\usepackage{amsmath,amssymb,amsthm}
\usepackage{booktabs}
\usepackage{longtable}
\usepackage{hyperref}
\usepackage{xcolor}

\newcommand{\code}[1]{\texttt{\small #1}}
\newcommand{\adj}[1]{\bar{#1}}
\newcommand{\dd}[2]{\frac{\partial #1}{\partial #2}}
\newcommand{\ddd}[3]{\frac{\partial^2 #1}{\partial #2 \, \partial #3}}
\newcommand{\gap}[1]{\textcolor{red}{\textbf{Gap:} #1}}

\title{Reverse-mode differentiation of the TF24 model:\\
the variables, the intermediaries, and the transposes}
\author{}
\date{2026-08-04}

\begin{document}
\maketitle

\begin{abstract}
This document states the mathematics of the reverse-mode trait gradient for the
TF24 size- and trait-structured model, with every intermediate quantity named and
mapped to the symbol that carries it in the code. Report 00 states the dependency
structure and walks the flow in prose; this states the derivatives. Its purpose is
that a reader can check a line of code against a line of algebra, in both
directions. Sections marked \gap{} are places where the algebra and the
implementation disagree today.
\end{abstract}

\tableofcontents

\section{What is being differentiated, and why the structure matters}

The model carries a density of individuals over a size coordinate, for each of
$S$ species, coupled to a shared light field and a shared soil water column. The
solver advances a state vector $y(t)$ by an embedded Runge--Kutta method, and
inserts new degrees of freedom as cohorts are introduced, so $\dim y$ grows during
a run.

We want
\begin{equation}
\nabla_\varphi \mathcal{C}
\qquad\text{where}\qquad
\mathcal{C} = \mathcal{C}\big(y(T), \varphi\big)
\label{eq:goal}
\end{equation}
is a scalar summary of the stand at time $T$ --- a census --- and
$\varphi \in \mathbb{R}^{P}$ collects the trait and physiological parameters. For
TF24, $P = 44$ per species.

Three properties of this problem set everything that follows.

\paragraph{Reverse mode is the only affordable direction.} A forward tangent costs
one solve per component of $\varphi$; the adjoint costs one solve plus one sweep
regardless of $P$. With $P = 44$ that is the difference between a feasible
computation and an infeasible one.

\paragraph{One layer of the model is an implicit function.} Each individual's
carbon gain is the value of a maximisation over its own stem water potential,
solved numerically. Differentiating through the solver is neither necessary nor
accurate; the envelope theorem and the implicit function theorem give the
derivatives in closed form, and Section~\ref{sec:leaf} does that.

\paragraph{The state grows.} An introduction widens $y$, so the adjoint is not a
single backward solve over a fixed space but a sequence of backward segments,
each narrower than the one after it.

\section{Notation, and the symbol that carries each quantity}
\label{sec:notation}

Every quantity below is either a state, a parameter, or an intermediary. The
third column is the name in the code; this table is the contract between the
algebra and the implementation.

\subsection{Indices}

\begin{longtable}{@{}lll@{}}
\toprule
Symbol & Meaning & Range \\
\midrule
$s$ & species & $1 \dots S$ \\
$k$ & cohort, ordered by the quadrature abscissa & $1 \dots N_s$ \\
$j$ & soil layer & $1 \dots L$ \\
$q$ & knot of the light field & $1 \dots K$ \\
$i$ & Runge--Kutta stage & $1 \dots 6$ \\
$n$ & accepted step of the solver & $1 \dots M$ \\
$p$ & component of $\varphi$ & $1 \dots P$ \\
\bottomrule
\end{longtable}

\subsection{State}

\begin{longtable}{@{}llp{0.42\textwidth}@{}}
\toprule
Symbol & Code & Meaning \\
\midrule
$h_k$ & \code{HEIGHT\_INDEX} & height of cohort $k$ \\
$\ell_k = \log n_k$ & \code{state(LOG\_DENSITY\_INDEX)} & log of the density carried by cohort $k$ \\
$b_k$ & \code{introduction\_time()} & the time cohort $k$ was introduced \\
$m^{\mathrm{hw}}_k, a^{\mathrm{hw}}_k$ & \code{MASS\_HEARTWOOD}, \code{AREA\_HEARTWOOD} & accumulated heartwood \\
$r_k$ & \code{NSC\_INDEX} & non-structural carbon reserve \\
$\theta_j$ & \code{soil\_moist} & water content of layer $j$ \\
$y$ & \code{Patch::ode\_state} & all of the above, species-major \\
\bottomrule
\end{longtable}

The quadrature abscissa is
\begin{equation}
x_k =
\begin{cases}
b_k & \text{birth-date coordinate} \\
-h_k & \text{height coordinate}
\end{cases}
\label{eq:abscissa}
\end{equation}
which is \code{Species::abscissa\_of}. It increases with $k$ in both cases. The
reverse-mode gradient is scoped to the birth-date coordinate, where $x_k = b_k$ is
fixed at birth and carries no derivative. This is the single most consequential
simplification in the design and Section~\ref{sec:reductions} shows why.

\subsection{Intermediaries: the ones that matter}

\begin{longtable}{@{}llp{0.40\textwidth}@{}}
\toprule
Symbol & Code & Meaning \\
\midrule
$A_k$ & \code{area\_leaf} & leaf area of cohort $k$, $\;(h_k/a_{l1})^{1/a_{l2}}$ \\
$E^{\mathrm{comp}}(z)$ & \code{compute\_competition} & competition experienced at height $z$ \\
$\Lambda_q, \Lambda'_q$ & knot values and slopes & the light field, as a spline \\
$p^\star_k$ & \code{opt\_psi\_stem\_}, \code{root\_collar\_psi\_} & the maximising stem water potential \\
$\Pi_k$ & \code{Leaf::profit\_} & carbon profit at the optimum \\
$U_{kj}$ & \code{soil\_consumption\_} & water cohort $k$ draws from layer $j$ \\
$\kappa_k$ & \code{leaf\_specific\_conductance\_max\_} & maximum leaf-specific conductance \\
$c^{\mathrm{i}}_k$ & \code{ci\_} & intercellular CO\textsubscript{2}, itself a root of a residual \\
$G(\cdot)$ & \code{build\_cumulative\_vulnerability\_integral} & cumulative vulnerability integral \\
$g_k$ & \code{rate(HEIGHT\_INDEX)} & height growth rate \\
$\mu_k$ & mortality rate & \\
\bottomrule
\end{longtable}

\section{The forward model as a composition}
\label{sec:forward}

Within one evaluation of the rates at one time, the model is a directed acyclic
graph. Writing it as a composition is what makes the transpose mechanical.

\begin{equation}
\varphi
\;\xrightarrow{\;\text{(a)}\;}\;
\text{strategy}
\;\xrightarrow{\;\text{(b)}\;}\;
\{A_k, \kappa_k\}
\;\xrightarrow{\;\text{(c)}\;}\;
\Lambda
\;\xrightarrow{\;\text{(d)}\;}\;
\{p^\star_k, \Pi_k, U_{kj}\}
\;\xrightarrow{\;\text{(e)}\;}\;
\dot y
\label{eq:composition}
\end{equation}

\begin{description}
\item[(a) \code{prepare\_strategy}] derives the strategy's dependent quantities
  from $\varphi$.
\item[(b) allometry] gives each cohort's size-dependent quantities from $h_k$ and
  $\varphi$.
\item[(c) the field reductions] aggregate over cohorts to give the shared
  environment. This is the only \emph{many-to-few} map in the chain, and therefore
  the only place a transpose scatters rather than gathers.
\item[(d) the individual] solves its own maximisation given the field. This is the
  only implicit step.
\item[(e) assembly] writes the rates back into $\dot y$.
\end{description}

The important structural fact is that (c) is a reduction and (d) is a
\emph{per-cohort independent} map. So the reverse pass is: transpose (e)
cohort-by-cohort, transpose (d) cohort-by-cohort, then transpose (c) once, which
scatters the field's adjoint back over every cohort.

\section{The discrete adjoint of the solver}
\label{sec:solver}

Let the solver advance $y_n \mapsto y_{n+1}$ by
\begin{align}
Y_i &= y_n + \Delta t_n \sum_{l<i} a_{il} k_l,
&
k_i &= f(Y_i, \varphi, t_n + c_i \Delta t_n),
\\
y_{n+1} &= y_n + \Delta t_n \sum_i \beta_i k_i .
\end{align}

Given the adjoint $\adj{y}_{n+1}$ of the state after the step, the step's
transpose is
\begin{equation}
\adj{k}_i = \Delta t_n \beta_i \adj{y}_{n+1}
  + \Delta t_n \sum_{l>i} a_{li} \adj{Y}_l,
\qquad
\adj{Y}_i = \left(\dd{f}{y}\Big|_{Y_i}\right)^{\!\!\top} \adj{k}_i ,
\label{eq:stageadj}
\end{equation}
swept for $i$ descending, and then
\begin{equation}
\adj{y}_n = \adj{y}_{n+1} + \sum_i \adj{Y}_i,
\qquad
\adj{\varphi} \mathrel{+}= \sum_i \left(\dd{f}{\varphi}\Big|_{Y_i}\right)^{\!\!\top}
\adj{k}_i .
\label{eq:paramaccum}
\end{equation}

Two consequences that the implementation must respect.

\paragraph{The step size is held.} $\Delta t_n$ was chosen by the adaptive
controller using an error estimate. Treating it as a differentiable function of
$\varphi$ would differentiate the controller, which is not the model. So
$\Delta t_n$ is a recorded constant on the reverse pass. This is why a recorded
trajectory must store its step sizes rather than recompute them.

\paragraph{$\adj\varphi$ accumulates over every step and every stage.} There are
$6M$ contributions, and \eqref{eq:paramaccum} is the only place they enter. In the
code that accumulator is \code{Patch::trait\_adjoint}.

\subsection{Introductions: the state changes dimension}

Let the run have widths $d_1 < d_2 < \dots < d_B$ over $B$ segments, with an
introduction between consecutive segments. The adjoint sweeps segment $B$ first.
At the boundary between segment $\sigma+1$ and segment $\sigma$, the newly
introduced components have no predecessor in $y$: their adjoints leave the state
and enter $\adj\varphi$ through the boundary condition
\begin{equation}
n_{\text{new}} = \frac{\text{birth\_rate} \cdot \text{pr\_estab}}{g},
\qquad
\ell_{\text{new}} = \log n_{\text{new}} ,
\label{eq:boundary}
\end{equation}
which reads $\varphi$ directly and is \code{Patch::introduction\_adjoint}. Then
$\adj y$ is narrowed to $d_\sigma$ and the sweep continues.

\emph{Requirement.} The narrowing must be applied to every column when several
functionals are swept together, and the segment list must cover every recorded
step with no gap. Nothing currently asserts the second.

\section{The recorded cohort step}
\label{sec:block}

The unit of the reverse pass is one evaluation of
\code{Individual::compute\_rates} for one cohort at one Runge--Kutta stage,
recorded at the active scalar type. Write its inputs as
\begin{equation}
u_k = \big(\underbrace{h_k, \ell_k, m^{\mathrm{hw}}_k, a^{\mathrm{hw}}_k, r_k}_{\text{own state}},\;
\underbrace{\Lambda, \Lambda'}_{\text{field, } 2K},\;
\underbrace{\theta_{1:L}}_{\text{soil}},\;
\underbrace{\varphi}_{P}\big)
\label{eq:blockin}
\end{equation}
and its outputs as
\begin{equation}
v_k = \big(\underbrace{\dot h_k, \dot m^{\mathrm{hw}}_k, \dot a^{\mathrm{hw}}_k, \dot r_k, \dot{(\text{fecundity})}, \mu_k}_{\text{rates}},\;
\underbrace{\dot \ell_k}_{\text{density rate}},\;
\underbrace{U_{k,1:L}}_{\text{consumption}}\big).
\label{eq:blockout}
\end{equation}

For TF24 with $K = 65$ and $L = 5$ this is $185$ inputs and $12$ outputs, and the
reverse pass forms $\big(\partial v_k / \partial u_k\big)^{\!\top} \adj v_k$.

\subsection{The density rate is where the coordinate choice bites}

On the birth-date coordinate,
\begin{equation}
\dot \ell_k = -\mu_k ,
\label{eq:densbirth}
\end{equation}
whereas on the height coordinate the density is compressed by the growth of the
size axis and
\begin{equation}
\dot \ell_k = -\mu_k - \dd{g}{h}\Big|_{h_k} .
\label{eq:densheight}
\end{equation}

Equation~\eqref{eq:densheight} is why the height coordinate is expensive to
differentiate: $\partial g/\partial h$ requires a \emph{second} solve of the whole
individual at a displaced height, so every recorded cohort step contains two
complete leaf maximisations, and the reverse pass must transpose both. Equation
\eqref{eq:densbirth} needs neither: $\mu_k$ is already an output of the same step.
Choosing the birth-date coordinate therefore halves the recorded work before any
optimisation is applied.

\section{The environment reductions and their transposes}
\label{sec:reductions}

\subsection{Light}

The competition profile at height $z$ sums over every cohort of every species:
\begin{equation}
E^{\mathrm{comp}}(z) = \sum_s \sum_k w_k \, n_k \, k_I^{(s)} \, A_k \,
\tilde{Q}\!\left(\frac{z}{h_k}\right),
\label{eq:light}
\end{equation}
where $\tilde{Q}(\nu) = (1 - \nu^{\eta})^2$ for $\nu \le 1$ and $0$ above, and
$w_k$ are the trapezium weights on the abscissa $x_k$ of \eqref{eq:abscissa}. The
field is then a spline through $K$ knots, so the recorded cohort step reads
$(\Lambda_q, \Lambda'_q)$ and not $E^{\mathrm{comp}}$ directly.

Given knot adjoints $\adj\Lambda_q$, the transpose of \eqref{eq:light} scatters:
\begin{align}
\adj\ell_k &\mathrel{+}= \sum_q \adj\Lambda_q \, w_k \, n_k \, k_I \, A_k \,
  \tilde{Q}(z_q/h_k),
\label{eq:lightadj_ell}\\
\adj h_k &\mathrel{+}= \sum_q \adj\Lambda_q \, w_k \, n_k \, k_I
  \left[ A_k' \, \tilde{Q}(z_q/h_k)
        - A_k \, \tilde{Q}'(z_q/h_k) \frac{z_q}{h_k^2} \right]
  \;+\; \underbrace{\sum_q \adj\Lambda_q \, n_k \, k_I A_k \tilde{Q}
        \dd{w_k}{h_k}}_{\text{\emph{zero on the birth-date coordinate}}},
\label{eq:lightadj_h}\\
\adj{k_I} &\mathrel{+}= \sum_q \adj\Lambda_q \sum_k w_k \, n_k \, A_k \,
  \tilde{Q}(z_q/h_k),
\label{eq:lightadj_kI}\\
\adj\eta &\mathrel{+}= \sum_q \adj\Lambda_q \sum_k w_k \, n_k \, k_I A_k \,
  \dd{\tilde{Q}}{\eta} .
\label{eq:lightadj_eta}
\end{align}

Two things to read off this.

The braced term in \eqref{eq:lightadj_h} is the weight-derivative term. It exists
only because the height coordinate makes the quadrature abscissa a function of the
state. On the birth-date coordinate $x_k = b_k$ is fixed at birth and passive, so
$\partial w_k/\partial h_k = 0$ and the term vanishes identically. Report 00
section 8 names this as the term a reader is most likely to forget; the coordinate
change removes the need to remember it.

Equations \eqref{eq:lightadj_kI} and \eqref{eq:lightadj_eta} are trait
contributions that arise \emph{inside the reduction}, not inside any cohort step.

\gap{The implementation has no path for \eqref{eq:lightadj_kI} or
\eqref{eq:lightadj_eta}. \code{Species::compute\_competition\_and\_slope\_adjoint}
returns \code{node\_size\_adjoints}, a structure with exactly three members ---
\code{area\_leaf}, \code{height}, \code{log\_density} --- and no parameter member,
and \code{Patch::trait\_adjoint} is written only from the cohort step and from the
introduction boundary. Since $k_I$ enters the model \emph{only} through
\eqref{eq:light}, its gradient is identically zero for every functional. The same
holds for $\eta$. This is a missing summation, not a hard derivative: both
right-hand sides are already computed as intermediate products inside the existing
transpose.}

\subsection{Water}

Total draw from layer $j$ aggregates the same way,
\begin{equation}
\mathcal{U}_j = \sum_s \sum_k w_k \, n_k \, U_{kj},
\label{eq:water}
\end{equation}
and the soil state responds through the retention curve
$\psi_j = \psi(\theta_j)$. Its transpose scatters $\adj{\mathcal{U}}_j$ back onto
$\adj U_{kj}$, $\adj \ell_k$ and, on the height coordinate, the weights.

\gap{\code{Species::consumption\_rate} now sorts the cohort grid when the
abscissa order inverts; \code{consumption\_rate\_adjoint} still transposes the
unsorted trapezium. On an inverted grid the two are not transposes of each other,
the gradient is finite, and nothing raises an error.}

\section{The individual's maximisation}
\label{sec:leaf}

This is the mathematical core. Each cohort chooses a stem water potential $p$ to
maximise profit,
\begin{equation}
p^\star = \arg\max_{p \in [p_a, p_b]} \Pi(p; u),
\qquad
\Pi = A\big(c^{\mathrm{i}}(p; u)\big) - \Theta(p; u),
\label{eq:argmax}
\end{equation}
with $A$ the assimilation and $\Theta$ the hydraulic cost. Everything downstream
--- profit, uptake, and hence growth --- is evaluated at $p^\star$, which is
computed numerically. We need derivatives of quantities at $p^\star$ with respect
to every input $u$, without differentiating the search.

\subsection{The interior optimum: profit is free}

Suppose $p^\star$ is interior, so stationarity holds:
\begin{equation}
g(p^\star, u) \;\equiv\; \dd{\Pi}{p}(p^\star, u) \;=\; 0 .
\label{eq:stationarity}
\end{equation}
Let $\Pi^\star(u) = \Pi(p^\star(u), u)$. Then
\begin{equation}
\dd{\Pi^\star}{u}
= \underbrace{\dd{\Pi}{p}}_{=\,0} \dd{p^\star}{u} + \dd{\Pi}{u}
= \dd{\Pi}{u} .
\label{eq:envelope}
\end{equation}
This is the envelope theorem, and it is the single largest economy in the design:
\textbf{the derivative of profit needs no sensitivity of the optimiser at all}.

\subsection{The argmax channel is rank one}

Other outputs are not stationary in $p$. Uptake $E(p^\star, u)$ is one. For those
we need $\partial p^\star/\partial u$. Differentiating \eqref{eq:stationarity},
\begin{equation}
\Pi_{pp} \, \dd{p^\star}{u} + \Pi_{pu} = 0
\qquad\Longrightarrow\qquad
\dd{p^\star}{u} = -\frac{\Pi_{pu}}{\Pi_{pp}} ,
\label{eq:argmaxsens}
\end{equation}
with $\Pi_{pp} = \partial^2\Pi/\partial p^2$ a scalar and
$\Pi_{pu} = \partial^2 \Pi/\partial p \, \partial u$ a row vector over the inputs.

In reverse mode we never form $\partial p^\star/\partial u$. Given an output
adjoint $\adj v$ on outputs $v = f(p^\star, u)$, define the scalar
\begin{equation}
s = \adj{v}^{\!\top} \dd{f}{p},
\qquad
m = -\frac{s}{\Pi_{pp}} ,
\label{eq:mu}
\end{equation}
and then
\begin{equation}
\adj u = \adj{v}^{\!\top}\dd{f}{u} \;+\; m \, \Pi_{pu}.
\label{eq:argmaxadj}
\end{equation}
Because $p$ is a \emph{scalar}, the whole optimiser channel is the outer product
$m \otimes \Pi_{pu}$ --- rank one. In the code $s$ is \code{s\_adjoint} and $m$ is
\code{mu}; the divisor $\Pi_{pp}$ is \code{dR\_dcollar\_at}.

\subsection{$\Pi_{pu}$ is the one genuinely new object}

Everything else in this section reuses derivatives the forward model already
needs. $\Pi_{pu}$ does not: it is a mixed second derivative of the profit, one
entry per input, and it must include the implicit-function term of the
$c^{\mathrm{i}}$ root-find (Section~\ref{sec:ci}). Report 00 calls it the single
new piece of code the design needs.

\gap{The implementation substitutes a two-sided finite difference of
$\partial \Pi/\partial p$ in $\varphi$ for the parameter half of $\Pi_{pu}$: for
each of the parameters that reach the operating point it perturbs the parameter and
re-evaluates. This costs $22$ of the $35$ residual evaluations per call, and its
conditioning has never been measured.}

\subsection{The pinned optimum: the envelope theorem does not apply}

If the maximiser sits at a bound, $p^\star = B(u)$ with $B$ either the
zero-uptake potential $p_a$ or the critical potential $p_b$, then
\eqref{eq:stationarity} is false: $\partial \Pi/\partial p \ne 0$ there. Instead
$p^\star$ follows the bound,
\begin{equation}
\dd{p^\star}{u} = \dd{B}{u},
\label{eq:pinnedsens}
\end{equation}
and the profit term reappears in the adjoint. With
\begin{equation}
w = \adj\Pi \, \dd{\Pi}{p} + s ,
\qquad
\adj u = \adj{v}^{\!\top}\dd{f}{u} + w \, \dd{B}{u},
\label{eq:pinnedadj}
\end{equation}
where $\partial B/\partial u$ is \code{Leaf::bound\_partials}. This is why the
pinned branch of the code carries a profit contribution that the interior branch
correctly omits: it is not an inconsistency, it is the envelope theorem failing at
a boundary.

The bound itself is implicit. $p_a$ is the collar potential at which uptake
vanishes, $E(p_a, u) = 0$, so
\begin{equation}
\dd{p_a}{u} = -\left(\dd{E}{p}\right)^{\!-1} \dd{E}{u} .
\label{eq:bounda}
\end{equation}

\subsection{The $c^{\mathrm{i}}$ root-find}
\label{sec:ci}

Intercellular CO\textsubscript{2} solves a supply-equals-demand residual
\begin{equation}
R(c^{\mathrm{i}}; p, u) = A(c^{\mathrm{i}})
  - \frac{g_c(p,u)\,(c^{\mathrm{a}} - c^{\mathrm{i}})}{P_{\mathrm{atm}}} = 0,
\label{eq:ci}
\end{equation}
bracketed on $[\Gamma^\star, c^{\mathrm{a}}]$. By the implicit function theorem
\begin{equation}
\dd{c^{\mathrm{i}}}{\bullet} = -\left(\dd{R}{c^{\mathrm{i}}}\right)^{\!-1}
\dd{R}{\bullet},
\label{eq:ciadj}
\end{equation}
and this term must be carried into every derivative that passes through
assimilation, including $\Pi_{pu}$.

The bracket is valid only when the endpoints straddle zero. Since
$R(c^{\mathrm{a}}) = A(c^{\mathrm{a}}) = A_{\max}$ and $g_c \propto E$, the
bracket fails when either $A_{\max} < 0$ or $E < 0$ --- both of which describe an
individual that is not producing. Those are exactly the states at which the
forward model substitutes $c^{\mathrm{i}} = \Gamma^\star$ with zero flux rather
than solving.

\gap{The forward path tests both conditions; the derivative path tests one; the
reverse path tests neither and calls the solver, which raises. The consequence is
not a bracket that is too narrow --- in the $A_{\max}$ case no root exists in the
domain at all --- but a missing guard.}

\subsection{The transport integral in closed form}
\label{sec:transport}

Flux from soil to collar integrates a stretched-exponential vulnerability curve,
\begin{equation}
G(m) = \int_0^m \exp\!\left(-\left(\frac{\sigma}{b}\right)^{c}\right)
\mathrm{d}\sigma
= \frac{b}{c}\, \gamma\!\left(\frac{1}{c},\, X\right),
\qquad
X = \left(\frac{m}{b}\right)^{c},
\label{eq:G}
\end{equation}
with $\gamma$ the lower incomplete gamma function, and $E = \kappa\,[\,G(\psi_{\mathrm{stem}}) - G(\psi_{\mathrm{up}})\,]$.

Writing $a = 1/c$ and using the everywhere-convergent series
\begin{equation}
\gamma(a, x) = x^{a} e^{-x} \sum_{n \ge 0}
\frac{x^{n}}{a(a+1)\cdots(a+n)}
\;\equiv\; x^{a} e^{-x} \, \Sigma(a,x),
\label{eq:gammaseries}
\end{equation}
both derivatives are available in closed form. The $x$ derivative is the
integrand,
\begin{equation}
\dd{\gamma}{x} = x^{a-1} e^{-x},
\label{eq:dgdx}
\end{equation}
and the $a$ derivative follows from $\partial_a$ of each term, since the $n$th
term is $x^n$ over $\prod_{l=0}^{n}(a+l)$:
\begin{equation}
\dd{\gamma}{a} = \log(x)\,\gamma(a,x)
  + x^{a} e^{-x} \sum_{n\ge 0}
    \left(-\,t_n \sum_{l=0}^{n} \frac{1}{a+l}\right),
\qquad
t_n = \frac{x^{n}}{a(a+1)\cdots(a+n)} .
\label{eq:dgda}
\end{equation}
So one loop with one extra accumulator gives value and both derivatives. Chaining
to the parameters with $\partial X/\partial b = -cX/b$,
$\partial X/\partial c = X \log(m/b)$ and $\partial a/\partial c = -1/c^2$:
\begin{align}
\dd{G}{m} &= \exp\!\left(-\left(m/b\right)^{c}\right),
\label{eq:dGdm}\\
\dd{G}{b} &= \frac{\gamma}{c} - X \,\dd{\gamma}{x},
\label{eq:dGdb}\\
\dd{G}{c} &= -\frac{b}{c^{2}}\gamma
  + \frac{b}{c}\left( X \log(m/b) \dd{\gamma}{x}
  - \frac{1}{c^{2}} \dd{\gamma}{a} \right).
\label{eq:dGdc}
\end{align}
Note that $b$ needs only \eqref{eq:dgdx}, which is elementary; only $c$ reaches
the series derivative \eqref{eq:dgda}.

\gap{The implementation tabulates $G$ on a grid of $100$ knots whose upper limit
is $b\,\log(100)^{1/c}$ and whose spacing is that limit over the resolution, under
a loop bound $\psi \le \psi_{\max}$. The knot \emph{count} therefore steps between
$100$ and $101$ under a relative perturbation of $10^{-6}$ in $b$ or $c$, and a
finite difference across that step is not a derivative. Measured errors are $47$,
$131$ and $10{,}245$ times the correct values. Equations
\eqref{eq:dGdm}--\eqref{eq:dGdc} remove the grid, not merely its cost.}

\section{Supplying the individual's derivatives by hand}
\label{sec:supplied}

The maximisation of Section~\ref{sec:leaf} is solved in double precision and is
not recorded on the tape. Its derivatives are supplied instead. For an output
$v$ whose value the solver has already produced, the recorded expression is
\begin{equation}
\tilde{v} = v + \sum_i \left(\dd{v}{u_i}\right)
\big(u_i - \mathrm{passive}(u_i)\big),
\label{eq:supplied}
\end{equation}
where $\mathrm{passive}(\cdot)$ strips the derivative and returns the value. Two
properties:
\begin{enumerate}
\item $\tilde{v} = v$ \emph{exactly}, in value, because every bracket is zero.
\item $\partial \tilde v/\partial u_i$ is the supplied $\partial v/\partial u_i$.
\end{enumerate}
So the tape carries the correct number and the hand-derived derivative, and the
same construction serves forward and reverse mode. This is
\code{TF24\_Strategy::graft} and \code{odelia::ode::supplied\_derivative}.

Property 1 has a sharp consequence for verification: \textbf{a finite difference
of the recorded step cannot see an error in a supplied derivative}, because
perturbing $u_i$ changes $v$ through the solver, not through the bracket. The
supplied derivatives must be checked against the individual's own algebra, never
against a difference of the step that consumes them.

\gap{Property 1 also means a non-finite supplied derivative corrupts the
\emph{value}: if $\partial v/\partial u_i$ is NaN then $\mathrm{NaN} \times 0$ is
NaN and $\tilde v$ is NaN. \code{Leaf::layer\_flux\_partials} returns all entries
NaN at a branch kink --- equal potentials, gravity balance, or a collar potential
within $10^{-8}$ of zero --- with no caller testing for it, so a kink makes the
gradient and the forward tangent NaN together. The plain double path is unaffected
because \eqref{eq:supplied} is only assembled at an active type.}

\section{The census functional and its two terms}
\label{sec:census}

A census is a quadrature over the size distribution of a per-individual metric,
\begin{equation}
\mathcal{C} = \sum_s \sum_k w_k \, n_k \, \mathfrak{m}(h_k, \varphi),
\label{eq:census}
\end{equation}
with $\mathfrak{m}$ one of leaf area, above-ground mass, or stem area. Because
$\mathfrak{m}$ reads $\varphi$ directly, the total derivative has two terms:
\begin{equation}
\dd{\mathcal{C}}{\varphi}
= \underbrace{\sum_s \sum_k w_k \, n_k \, \dd{\mathfrak{m}}{\varphi}}
  _{\text{direct, at fixed state}}
\;+\;
\underbrace{\left(\dd{\mathcal{C}}{y}\right)^{\!\top}
  \dd{y}{\varphi}}_{\text{through the trajectory, by Section~\ref{sec:solver}}} .
\label{eq:censustwo}
\end{equation}

The second term is what the adjoint sweep computes, seeded with
$\adj y(T) = \partial \mathcal{C}/\partial y$. The first is not a sensitivity of
the state at all and no sweep produces it.

For the concrete metrics: $\mathfrak{m} = A(h)$ reads $a_{l1}, a_{l2}$;
$\mathfrak{m} = m^{\text{leaf}} = A(h)\,\mathrm{lma}$ reads $\mathrm{lma}$;
stem area reads $\theta$ and $a_{b1}$ through sapwood and bark.

\gap{\code{SCM::census\_state\_adjoint} registers only $y$ as an input, so the
direct term of \eqref{eq:censustwo} is absent. Note also that the seed's support
is wider than the metric algebra suggests: the recording calls
\code{set\_ode\_state}, which rebuilds the boundary node \eqref{eq:boundary},
whose density is a full physiology evaluation through the light field. So most of
$\varphi$ takes a non-zero direct term, not only the parameters that appear in
$\mathfrak{m}$.}

\subsection{Several functionals share one recording}

For $F$ functionals the record is common and only the seed differs. Recording
once and sweeping $F$ times costs one record plus $F$ sweeps, against $F$ records
and $F$ sweeps; the record dominates.

\gap{The implementation records $F$ times, and worse: it constructs the copy at
the active type \emph{once, outside} the loop, while each sweep begins by clearing
the tape. Clearing returns the derivative-slot counter to zero, so every active
value that outlives a sweep refers to a slot that now belongs to something else.
The first functional is correct and every later one reads unrelated storage.
Measured: rows agree with an independent reference for the first metric and, for
the later ones, have the wrong sign, a magnitude wrong by $180$ times, and $33$ of
$52$ columns exactly zero. This is the rule stated in the project's own method
document and obeyed at the two other sites that build such copies.}

\section{Where the trait gradient is assembled}

Collecting the paths by which $\varphi$ reaches $\mathcal{C}$:
\begin{equation}
\adj\varphi =
\underbrace{\sum_s\sum_k w_k n_k \dd{\mathfrak{m}}{\varphi}}_{\text{census, direct}}
+ \underbrace{\sum_n \sum_i \left(\dd{f}{\varphi}\right)^{\!\top}\!\adj k_i}
  _{\text{cohort steps, all stages}}
+ \underbrace{\sum_{\text{intro}} \dd{\ell_{\text{new}}}{\varphi}\,\adj\ell_{\text{new}}}
  _{\text{boundary}}
+ \underbrace{\sum_q \adj\Lambda_q \dd{E^{\mathrm{comp}}}{\varphi}}
  _{\text{field reduction}} .
\label{eq:assemble}
\end{equation}

The four terms are the four places $\varphi$ is read. The implementation has the
second and the third; the first is Section~\ref{sec:census}'s gap and the fourth
is Section~\ref{sec:reductions}'s.

\subsection{One term is absent from every path}

\code{rebind\_from} constructs the copy at the active type by copying the
strategy, but \code{prepare\_strategy} and the seed height refuse an active
scalar, so $h_0$, $A_0$ and $\eta_c$ enter as values. Hence
\begin{equation}
\dd{h_0}{\varphi} = 0
\label{eq:heightseed}
\end{equation}
is imposed on both AD paths, forward and reverse.

\gap{Equation \eqref{eq:heightseed} is false. The seed height solves an implicit
condition on the strategy and does depend on $\varphi$. Measured as about $3$ per
cent for $\mathrm{lma}$. \textbf{No instrument in the present design can referee a
fix}: the forward tangent imposes the same equation, and a re-run finite difference
cannot referee at production because a relative step of $2\times 10^{-7}$ in
$\mathrm{lma}$ moves a mature stand between alive and identically zero. Closing it
needs an implicit-value treatment of the seed height \emph{and} a new reference.}

\section{Summary of the gaps, in dependency order}

\begin{longtable}{@{}rp{0.52\textwidth}l@{}}
\toprule
& Gap & Section \\
\midrule
1 & Copy at the active type outlives a cleared tape; every functional after the
    first reads unrelated storage & \ref{sec:census} \\
2 & Field reduction has no parameter accumulator, so $k_I$ and $\eta$ are
    identically zero & \ref{sec:reductions} \\
3 & Census direct term not registered & \ref{sec:census} \\
4 & No guard on the $c^{\mathrm{i}}$ bracket at a non-producing individual &
    \ref{sec:ci} \\
5 & NaN supplied derivative at a branch kink corrupts the value &
    \ref{sec:supplied} \\
6 & Transport derivatives differenced across a grid whose knot count moves &
    \ref{sec:transport} \\
7 & $\Pi_{pu}$'s parameter half differenced, conditioning unmeasured &
    \ref{sec:leaf} \\
8 & Consumption adjoint not the transpose of a sorted forward grid &
    \ref{sec:reductions} \\
9 & $\mathrm{d}h_0/\mathrm{d}\varphi$ imposed zero on both AD paths; no
    available reference & \ref{eq:heightseed} \\
\bottomrule
\end{longtable}

Items 1 to 3 are missing summations or misplaced constructions and each is a small
change. Items 4 and 5 are missing guards. Items 6 and 7 replace a difference with
algebra this document gives in closed form. Item 8 is a transpose that stopped
matching its forward function. Item 9 is the only one that needs a new instrument
before it can be closed, and it is therefore the only one whose scope is genuinely
open.

\end{document}
